% Part: proof-theory
% Chapter: normalization
% Section: normalization-thm

\documentclass[../../../include/open-logic-section]{subfiles}

\begin{document}

\olfileid{pt}{nor}{thm}

\olsection{સામાન્યીકરણ પ્રમેય}

!!^a{proof}માં કટ-ખંડ
ન હોય તો તે
\emph{સામાન્ય} સ્વરૂપમાં
છે. સ્પષ્ટ છે કે
આ ત્યારે અને ત્યારે
જ બને જ્યારે તેના
પર ક્રમપરિવર્તન-રૂપાંતર
કે ઘટાડા-રૂપાંતર,
બેમાંથી એકેય લગાડી
ન શકાય.
!!a{proof}નું
\emph{સામાન્યીકરણ} કરવું
એટલે કટ ન રહે
ત્યાં સુધી આ બંને
પ્રકારનાં રૂપાંતર
લગાવી તેને સામાન્ય
સ્વરૂપની સાબિતીમાં
ફેરવવી. આવું હંમેશાં
કરી શકાય છે, એ
\emph{સામાન્યીકરણ પ્રમેય}
છે.

\begin{thm}\ollabel{thm:normalization}
દરેક \Log{N1i}-!!{proof}
$\delta$નું સામાન્યીકરણ
થાય છે; એટલે તેને
કટ વિનાની
!!a{proof}~$\delta'$માં
ફેરવી શકાય છે.
\end{thm}

\begin{proof}
  $\delta$ની કટ-કક્ષા
  અને કટ-લંબાઈ પર
  આગમન કરીએ.
  આગમન-પરિકલ્પના એવી
  છે કે ઓછી કટ-કક્ષા
  ધરાવતી, અથવા એ
  જ કટ-કક્ષા સાથે
  ઓછી કટ-લંબાઈ ધરાવતી,
  દરેક !!{proof}નું
  સામાન્યીકરણ થાય છે.

  કટ-લંબાઈ~$0$
  હોય તો કટ જ
  નથી અને~$\delta$
  પહેલેથી સામાન્ય
  સ્વરૂપમાં છે.
  એટલે માનો કે
  $\delta$ની કટ-કક્ષા
  અને કટ-લંબાઈ
  બંને~$> 0$ છે.
  સૌથી ઉપરનો અને
  સૌથી જમણો મહત્તમ
  કટ-ખંડ પસંદ કરો.
  તેની લંબાઈ~$>1$
  હોય તો
  ક્રમપરિવર્તન-રૂપાંતર
  લગાવો. લંબાઈ~$1$
  હોય તો અનુરૂપ
  ઘટાડા-રૂપાંતર
  લગાવો. આમ મળતી
  નવી !!{proof}ની
  કટ-કક્ષા સમાન
  અને કટ-લંબાઈ ઓછી
  હોય છે, અથવા
  કટ-કક્ષા જ ઓછી
  હોય છે; તેથી
  આગમન-પરિકલ્પના
  લાગુ પડે છે.
\end{proof}

ઉપરની સાબિતીની રીતમાં
હંમેશાં સૌથી ઉપરના,
સૌથી જમણા મહત્તમ
કટને પહેલાં હાથ
ધરીએ છીએ. તેથી
!!a{proof} પર રૂપાંતરો
ચોક્કસ ક્રમે જ
લગાવવાં પડે એવું
લાગે, જેથી સામાન્ય
!!{proof} મળે.
આશ્ચર્યની વાત એ
છે કે ક્રમ ખરેખર
મહત્ત્વનો નથી.
ક્રમપરિવર્તન અને
ઘટાડા-રૂપાંતરો કોઈપણ
ક્રમે લગાવતાં છેવટે
સામાન્ય સ્વરૂપની
સાબિતી મળે છે.

ખંડ અને કટની
વ્યાખ્યાઓ~\Log{N2i}
માટે સહેલાઈથી આપી
શકાય છે; અને
\Log{N1i} માટેનાં
ક્રમપરિવર્તન અને
ઘટાડા-રૂપાંતરો સીધાં
\Log{N2i}માં રૂપાંતરિત
થાય છે. તેથી
સામાન્યીકરણ પ્રમેય
\Log{N2i} માટે પણ
લાગુ પડે છે.

\begin{thm}\ollabel{thm:normalization-N2i}
દરેક \Log{N2i}-!!{proof}
$\delta$નું સામાન્યીકરણ
થાય છે; એટલે તેને
કટ વિનાની
!!a{proof}~$\delta'$માં
ફેરવી શકાય છે.
\end{thm}

\end{document}
